\chapter{Cząstka w nieskończonej studni potencjału}

\begin{summarybox}
Pierwszy model laboratoryjny: wyprowadzenie poziomów energetycznych, funkcji
własnych, normalizacji oraz wykresów \(\psi_n(x)\) i \(|\psi_n(x)|^2\).
Nieskończona studnia potencjału jest najprostszym miejscem, w którym widać,
że kwantowanie energii nie jest dodatkiem interpretacyjnym, lecz skutkiem
równania różniczkowego i warunków brzegowych.
\end{summarybox}

\section{Sytuacja fizyczna: cząstka uwięziona w pudełku}

Rozważamy cząstkę o masie \(m\), która może poruszać się tylko w pewnym skończonym przedziale. Poza tym przedziałem potencjał jest nieskończenie duży. Fizycznie oznacza to idealne, nieprzenikalne ściany. Matematycznie oznacza to, że funkcja falowa musi znikać na ścianach i poza nimi.

Najczęściej używamy modelu
\[
V(x)=
\begin{cases}
0, & 0<x<a,\\
\infty, & x\leq 0\ \text{lub}\ x\geq a.
\end{cases}
\]
Wewnątrz studni cząstka jest swobodna, bo \(V(x)=0\). Nie znaczy to jednak, że jej energia może być dowolna. Ograniczenie przestrzenne wymusza szczególne kształty funkcji falowej, a te kształty wymuszają dyskretne energie.

Ten model jest idealizacją. Prawdziwe ściany potencjału nie są nieskończone. Mimo to nieskończona studnia jest bardzo ważna: uczy, jak warunki brzegowe wybierają dozwolone rozwiązania, i daje prosty punkt odniesienia dla kropek kwantowych, nanostruktur oraz modeli elektronów w ograniczonych obszarach.

\section{Dwie konwencje studni: \texorpdfstring{\(x\in[0,a]\)}{x in [0,a]} oraz \texorpdfstring{\(x\in[-a/2,a/2]\)}{x in [-a/2,a/2]}}

Istnieją dwie popularne konwencje położenia studni. Pierwsza używa przedziału
\[
0\leq x\leq a.
\]
Druga używa przedziału symetrycznego
\[
-\frac{a}{2}\leq x\leq \frac{a}{2}.
\]
Obie opisują tę samą fizykę, jeśli szerokość studni wynosi \(a\). Różnią się tylko wyborem początku układu współrzędnych.

W przedziale \([0,a]\) naturalne funkcje własne mają postać sinusów:
\[
\psi_n(x)=\sqrt{\frac{2}{a}}\sin\left(\frac{n\pi x}{a}\right),
\qquad n=1,2,3,\ldots
\]
W przedziale symetrycznym wygodniej mówić o parzystości: część stanów jest parzysta, a część nieparzysta. Pojawiają się wtedy funkcje typu cosinusowego i sinusowego.

W tym rozdziale najpierw użyjemy konwencji \([0,a]\), bo prowadzi do najkrótszego rachunku. Potem wrócimy do konwencji symetrycznej, aby zobaczyć, jak widać parzystość stanów.

\begin{remarkbox}
Zmiana konwencji położenia studni nie zmienia energii. Energia zależy od szerokości \(a\), a nie od tego, gdzie ustawimy początek osi \(x\).
\end{remarkbox}

\section{Stacjonarne równanie Schrödingera}

Wewnątrz studni potencjał jest równy zero, więc stacjonarne równanie Schrödingera ma postać
\begin{equation}
-\frac{\hbar^2}{2m}\frac{d^2\psi}{dx^2}=E\psi.
\end{equation}
Przekształcamy je do postaci
\begin{equation}
\frac{d^2\psi}{dx^2}+k^2\psi=0,
\end{equation}
gdzie
\begin{equation}
k^2=\frac{2mE}{\hbar^2}.
\end{equation}
Ogólne rozwiązanie wewnątrz studni jest kombinacją sinusa i cosinusa:
\begin{equation}
\psi(x)=A\sin(kx)+B\cos(kx).
\end{equation}
Stałe \(A\), \(B\) oraz dozwolone wartości \(k\) nie są jeszcze znane. Wyznaczą je warunki brzegowe.

Na tym etapie rachunek wygląda jak klasyczne równanie falowe. Różnica pojawia się przy interpretacji i przy warunkach brzegowych: tylko niektóre długości fali mieszczą się w pudełku w sposób zgodny z zerowaniem funkcji falowej na ścianach.

\section{Warunki brzegowe i kwantowanie energii}

Dla studni \([0,a]\) funkcja falowa musi spełniać
\begin{equation}
\psi(0)=0,
\qquad
\psi(a)=0.
\end{equation}
Pierwszy warunek daje
\[
\psi(0)=A\sin 0+B\cos 0=B=0.
\]
Zostaje więc
\[
\psi(x)=A\sin(kx).
\]
Drugi warunek daje
\[
\psi(a)=A\sin(ka)=0.
\]
Nie interesuje nas rozwiązanie \(A=0\), bo oznaczałoby brak cząstki. Musi więc zachodzić
\begin{equation}
\sin(ka)=0.
\end{equation}
Stąd
\begin{equation}
ka=n\pi,
\qquad n=1,2,3,\ldots
\end{equation}
czyli
\begin{equation}
k_n=\frac{n\pi}{a}.
\end{equation}
Ponieważ \(E=\hbar^2 k^2/(2m)\), otrzymujemy dozwolone energie
\begin{equation}
E_n=\frac{n^2\pi^2\hbar^2}{2ma^2},
\qquad n=1,2,3,\ldots
\end{equation}

To jest pierwszy zasadniczy wynik: energia cząstki w pudełku jest skwantowana. Nie wstawiliśmy kwantowania ręcznie. Ono wyszło z równania Schrödingera i warunków brzegowych.

\begin{summarybox}
Warunki brzegowe \(\psi(0)=0\) i \(\psi(a)=0\) wymuszają \(k_n=n\pi/a\). Dlatego energie są dyskretne i rosną jak \(n^2\).
\end{summarybox}

\section{Funkcje własne i ich normalizacja}

Dla każdego \(n\) funkcja własna ma postać
\[
\psi_n(x)=A_n\sin\left(\frac{n\pi x}{a}\right).
\]
Stałą \(A_n\) wyznaczamy z warunku normalizacji:
\begin{equation}
\int_0^a |\psi_n(x)|^2\,dx=1.
\end{equation}
Dostajemy
\[
1=|A_n|^2\int_0^a \sin^2\left(\frac{n\pi x}{a}\right)dx.
\]
Dla całkowitego \(n\) zachodzi
\[
\int_0^a \sin^2\left(\frac{n\pi x}{a}\right)dx=\frac{a}{2}.
\]
Jeżeli wybieramy \(A_n\) rzeczywiste i dodatnie, to
\begin{equation}
A_n=\sqrt{\frac{2}{a}}.
\end{equation}
Ostatecznie
\begin{equation}
\psi_n(x)=\sqrt{\frac{2}{a}}\sin\left(\frac{n\pi x}{a}\right),
\qquad 0<x<a.
\end{equation}
Poza studnią funkcja falowa jest równa zero.

Funkcje te są ortonormalne:
\begin{equation}
\int_0^a \psi_n(x)\psi_m(x)\,dx=\delta_{nm}.
\end{equation}
Oznacza to, że różne stany energetyczne są wzajemnie ortogonalne i mogą tworzyć bazę do rozwijania bardziej ogólnych stanów w pudełku.

\section{Parzystość funkcji falowych}

W konwencji symetrycznej \([-a/2,a/2]\) dobrze widać parzystość funkcji falowych. Warunki brzegowe mają wtedy postać
\[
\psi\left(-\frac{a}{2}\right)=0,
\qquad
\psi\left(\frac{a}{2}\right)=0.
\]
Rozwiązania można podzielić na parzyste i nieparzyste.

Stany parzyste mają postać cosinusową:
\begin{equation}
\psi(x)=C\cos(kx),
\end{equation}
a stany nieparzyste mają postać sinusową:
\begin{equation}
\psi(x)=C\sin(kx).
\end{equation}
Parzystość oznacza odpowiednio
\[
\psi(-x)=\psi(x)
\]
albo
\[
\psi(-x)=-\psi(x).
\]

Fizyka pozostaje taka sama jak w konwencji \([0,a]\). Zmienia się tylko wygoda opisu. Jeżeli potencjał jest symetryczny względem zera, wybór przedziału symetrycznego pozwala natychmiast klasyfikować stany według parzystości.

\begin{remarkbox}
Parzystość jest prostym przykładem symetrii. Jeżeli potencjał jest symetryczny, to rozwiązania równania Schrödingera można wybrać jako parzyste albo nieparzyste.
\end{remarkbox}

\section{Poziomy energetyczne w eV}

W zastosowaniach wygodnie wyrażać energię w elektronowoltach. Wzór
\[
E_n=\frac{n^2\pi^2\hbar^2}{2ma^2}
\]
pokazuje trzy ważne zależności. Po pierwsze, energia rośnie jak \(n^2\). Po drugie, maleje jak \(a^{-2}\). Po trzecie, maleje wraz ze wzrostem masy cząstki.

Dla elektronu w studni o szerokości rzędu nanometra energie mogą być rzędu elektronowoltów. To tłumaczy, dlaczego efekty kwantowego ograniczenia są tak ważne w nanostrukturach. Jeśli szerokość pudełka zwiększymy dziesięć razy, energia spadnie sto razy.

W Mathematice można szybko policzyć poziomy energetyczne na przykład tak:
\begin{verbatim}
hbar = 1.054571817*10^-34;
me = 9.1093837015*10^-31;
eV = 1.602176634*10^-19;
a = 1*10^-9;

E[n_] := n^2 Pi^2 hbar^2/(2 me a^2)/eV
Table[{n, N[E[n]]}, {n, 1, 5}]
\end{verbatim}

Taki rachunek należy zawsze sprawdzić wymiarowo. Jeśli \(a\) jest podane w metrach, masa w kilogramach, a \(\hbar\) w jednostkach SI, to wynik przed podzieleniem przez \(eV\) jest w dżulach.

\section{Wizualizacja \texorpdfstring{\(\psi_n(x)\)}{psi_n(x)} oraz \texorpdfstring{\(|\psi_n(x)|^2\)}{|psi_n(x)|^2}}

Wykres funkcji falowej pokazuje znak i węzły stanu. Wykres \(|\psi_n(x)|^2\) pokazuje gęstość prawdopodobieństwa. Dla \(n=1\) funkcja falowa nie ma węzła wewnątrz studni. Dla \(n=2\) ma jeden węzeł wewnętrzny, dla \(n=3\) dwa węzły i tak dalej.

Przykładowy kod w Mathematice:
\begin{verbatim}
a = 1;
psi[n_, x_] := Sqrt[2/a] Sin[n Pi x/a]

Plot[Evaluate[Table[psi[n, x], {n, 1, 4}]], {x, 0, a},
 PlotLegends -> Table["n=" <> ToString[n], {n, 1, 4}],
 AxesLabel -> {"x", "psi_n(x)"}]

Plot[Evaluate[Table[psi[n, x]^2, {n, 1, 4}]], {x, 0, a},
 PlotLegends -> Table["n=" <> ToString[n], {n, 1, 4}],
 AxesLabel -> {"x", "|psi_n(x)|^2"}]
\end{verbatim}

W interpretacji fizycznej trzeba uważać: ujemna wartość \(\psi_n(x)\) nie oznacza ujemnego prawdopodobieństwa. Prawdopodobieństwo jest związane z modułem kwadratowym. Znak funkcji falowej ma jednak znaczenie przy superpozycjach i interferencji.

\section{Zależność energii od szerokości studni}

Najbardziej charakterystyczna zależność ma postać
\begin{equation}
E_n\propto \frac{1}{a^2}.
\end{equation}
Im ciaśniej ograniczamy cząstkę, tym większa jest jej minimalna energia. To jest konkretna realizacja intuicji związanej z zasadą nieoznaczoności: mała niepewność położenia wymusza większą niepewność pędu, a więc większą energię kinetyczną.

Dla stanu podstawowego
\begin{equation}
E_1=\frac{\pi^2\hbar^2}{2ma^2}.
\end{equation}
Jeżeli \(a\) zmniejszymy dwukrotnie, energia wzrośnie czterokrotnie. Jeśli \(a\) zmniejszymy dziesięciokrotnie, energia wzrośnie stukrotnie.

To zachowanie jest jednym z powodów, dla których nanoskala jest fizycznie szczególna. Zmiana rozmiaru układu może przesuwać poziomy energetyczne na tyle silnie, że zmienia własności optyczne i transportowe materiału.

\section{Masa efektywna i przykład GaAs}

W półprzewodnikach elektron w krysztale często zachowuje się tak, jakby miał inną masę niż elektron swobodny. Nazywamy ją masą efektywną i oznaczamy \(m^*\). W prostym modelu studni wzór na energię przyjmuje wtedy postać
\begin{equation}
E_n=\frac{n^2\pi^2\hbar^2}{2m^*a^2}.
\end{equation}

Dla GaAs masa efektywna elektronu w paśmie przewodnictwa jest znacznie mniejsza od masy elektronu swobodnego, często używa się przybliżenia
\[
m^*\approx 0.067m_e.
\]
Ponieważ energia jest odwrotnie proporcjonalna do masy, mniejsza masa efektywna oznacza większe odstępy między poziomami energetycznymi dla tej samej szerokości studni.

To pokazuje, dlaczego ten sam model matematyczny może opisywać różne układy fizyczne: zmieniamy parametry, takie jak masa efektywna i szerokość studni, a otrzymujemy inny zakres energii. Właśnie dlatego nieskończona studnia potencjału jest podstawowym modelem roboczym dla struktur kwantowych.

\begin{summarybox}
Nieskończona studnia potencjału pokazuje w najczystszej postaci trzy idee: warunki brzegowe wymuszają kwantowanie, funkcje własne tworzą bazę stanów, a energia rośnie jak \(n^2/a^2\). Ten prosty model będzie punktem odniesienia dla skończonej studni, tunelowania i struktur półprzewodnikowych.
\end{summarybox}
